In conventional block copolymers, spontaneous curvature is usually tuned by prescribed molecular parameters, such as the Kuhn length, component volume fraction, or grafting number\cite{Xu2024Architectural}.
In contrast, L--S TT copolymers can locally adjust the preferred topology $f_{\mathcal{R}}^\ast$ through sliding, because AB linear diblock and A$_2$B star topologies possess different curvature preferences.
As a weakly curved three-dimensional network phase whose interfacial curvature is closest to that of the lamellar phase, the $\mathrm{O}^{70}$ phase provides a useful model for examining how topology transformability affects interfacial adaptation.
We therefore use the $\mathrm{O}^{70}$ phases at low and high $\bar{\phi}_{\mathrm A}$ in the low-$N\chi$ region to examine how the sliding degree of freedom affects microphase separation and responds to interfacial curvature.


As shown in Fig.~\ref{fig:PhaseDiagram}, increasing $\phi_\mathrm{a}$ shifts the $\mathrm{O}^{70}$ phase region toward lower $\bar{\phi}_\mathrm{A}$ relative to the linear AB diblock reference in both the A-core-forming regime ($\bar{\phi}_\mathrm{A}<0.5$) and the A-shell-forming regime ($\bar{\phi}_\mathrm{A}>0.5$).
This trend originates from the reduction of the radius of gyration $R_g$ of the A block caused by the sliding side chain, which decreases the spontaneous curvature associated with the A component.
However, the shifts along the $N\chi$ direction differ between the A-core-forming and A-shell-forming regimes.
Beyond the slight increase in $N\chi_{\mathrm{ODT}}^{\mathrm{Slide}}$ relative to the linear AB diblock reference, the $\mathrm{O}^{70}$ phase region shifts to higher $N\chi$ in the A-core-forming regime but to lower $N\chi$ in the A-shell-forming regime.
This asymmetry reflects the opposite curvature requirements of the A component in the two regimes: positive spontaneous curvature when A forms the core and negative spontaneous curvature when A forms the shell regions.
Because the sliding side chain reduces the spontaneous curvature of the A block, the statistical weights of different $f$ respond differently to the curvature environments in the two regimes.
The A-core-forming regime is associated with a larger contribution from more negative-curvature topologies, whereas the A-shell-forming regime favors nearly linear configurations that are more compatible with its local geometry.
This curvature-conditioned topology response is consistent with the asymmetric displacement of the $\mathrm{O}^{70}$ phase region in the phase diagram.


\begin{figure}[!htb]
    \centering
    \includegraphics[]{./Figure/O70-Pf.png}
    \caption{
        (a,b) Three-dimensional views of the $\mathrm{O}^{70}$ phase formed by L--S TT copolymers at $\bar{\phi}_{\mathrm A}=0.42$ and $0.55$, respectively.
        Red--yellow regions denote the A-component density.
        (c,d) Isodensity-layer topology distributions $P_\Omega(f)$ at $\phi_{\mathrm A}^{\mathrm{iso}}=0.75$ as functions of the sliding side-chain fraction $\phi_{\mathrm a}$ for the two compositions in panels (a) and (b), respectively.
        Heat maps show $P_\Omega(f)$, contours indicate numerical isovalues, and the red and blue lines denote the locations of the maximum and minimum of the distribution, $f_\Omega^\ast$ and $f_\Omega^-$, with symbols showing the corresponding numerical results.
        (e,f) Isodensity-layer topology distributions $P_\Omega(f)$ as functions of the sampled density level $\phi_{\mathrm A}^{\mathrm{iso}}$ at fixed $\phi_{\mathrm a}=0.50$, for $\bar{\phi}_{\mathrm A}=0.42$ and $0.55$, respectively, with the same notation as in panels (c,d).
        (g,h) Schematic illustrations of the topology redistribution of sliding chains across the interfacial region for $\bar{\phi}_{\mathrm A}=0.42$ and $0.55$, respectively.
    }
    \label{fig:O70Topology}
\end{figure}


To examine the topology response in the O$^{70}$ phase, we compare the Core ($\bar{\phi}_{\mathrm A}=0.42$) and Shell ($\bar{\phi}_{\mathrm A}=0.55$) states at $N\chi=13.0$, both with Fddd symmetry (Fig.~\ref{fig:O70Topology}a,b).
Fig.~\ref{fig:O70Topology}c,d shows their isodensity layer topology distributions $P_\Omega(f)$ at $\phi_{\mathrm A}^{\mathrm{iso}}=0.75$ as $\phi_{\mathrm a}$ varies.
At this A-rich density level, $P_\Omega(f)$ remains single peaked over the examined $\phi_{\mathrm a}$ range, with the red line tracing the preferred topology $f_\Omega^\ast$.
At small $\phi_{\mathrm a}$, $f_\Omega^\ast$ remains near $0.3$, corresponding to a typical ``Y''-shaped $\mathrm{A}_2\mathrm{AB}$ topology, and then decreases toward the AB-like limit as $\phi_{\mathrm a}$ increases.
This ordered-phase topology preference is consistent with the RPA prediction near the ODT and motivates the following analysis of its spatial relation to interfacial geometry.


To resolve the spatial dependence of topology selection, we examine $P_\Omega(f)$ across isodensity layers characterized by $\phi_{\mathrm A}^{\mathrm{iso}}$ at fixed $\phi_{\mathrm a}=0.50$.
As shown in Fig.~\ref{fig:O70Topology}e,f, the sampled layers span the A-rich, interfacial, and B-rich regions as $\phi_{\mathrm A}^{\mathrm{iso}}$ decreases from $0.90$ to $0.10$.
In the A-rich region, $P_\Omega(f)$ is single peaked, with its maximum at $f_\Omega^\ast$, whereas toward the B-rich region it develops a U-shaped profile with enhanced weights near the AB-like ($f=0$) and A$_2$B-like ($f=1$) limits and a local minimum $f_\Omega^-$ between them.
For $\phi_{\mathrm{A}}^{\mathrm{iso}} \lesssim 0.40$, the endpoint contributions become increasingly pronounced while intermediate sliding configurations remain populated.
This layer dependence reveals a spatial redistribution of topology weights across the interface, as illustrated schematically in Fig.~\ref{fig:O70Topology}g,h.
The A-segment population in the A-rich region is preferentially contributed by A$_2$AB-like configurations, whereas toward the B-rich region the relative contributions shift toward the AB-like and A$_2$B-like limits.
This redistribution is consistent with their different natural extension lengths: the more compact A$_2$B-like configurations preferentially contribute near the interface, whereas the more extended AB-like configurations can reach farther toward the A-rich region while their B blocks remain closer to the B-rich domain interior.
The enhanced A$_2$AB-like contribution in the A-rich region reflects the participation of its two A branches in filling that domain.
The sliding degree of freedom therefore redistributes statistical weight within a single continuous topology ensemble across different spatial regions, rather than producing distinct topological populations; this interpretation is supported by the weak spatial variation of $\bar f_{\mathcal R}$ and the small $D_{\mathcal R}$ discussed in Sec.~\ntxref{SI-subsec:TopologyStatisticsSpatial} of the Supporting Information.



% The response of $P_\Omega(f)$ to $\phi_{\mathrm A}^{\mathrm{iso}}$ further demonstrates the interfacial-curvature adaptability unique to L--S TT copolymers.
% This adaptability manifests differently in the two $\mathrm{O}^{70}$ phases because the A component plays distinct structural roles in the A-core-forming and A-shell-forming regimes.
% For $\bar{\phi}_{\mathrm A}=0.42$, where A forms the network core, $f_\Omega^\ast$ jumps from the A$_2$B endpoint to the AB linear diblock endpoint near $\phi_{\mathrm A}^{\mathrm{iso}}\approx0.23$, stays close to zero over an intermediate range, and then increases toward an A$_2$AB-like topology at larger $\phi_{\mathrm A}^{\mathrm{iso}}$.
% This behavior can be attributed to the larger spatial volume occupied by the two A branches in the A$_2$B topology, which provides a lower, or even negative, A-side spontaneous curvature.
% Such a topology can compensate for the increased positive curvature associated with $\bar{\phi}_{\mathrm A}<0.5$ at the interface, thereby satisfying the curvature requirement of the $\mathrm{O}^{70}$ phase.
% In the A-core-forming geometry, the outer region of the A core requires additional A segments to maintain density continuity.
% As the interfacial stretching penalty increases, the preferred configuration shifts from an AB linear diblock topology toward a Y-shaped A$_2$AB topology, filling the interfacial region more efficiently and lowering the free energy.
% By contrast, for $\bar{\phi}_{\mathrm A}=0.55$, the low-$\phi_{\mathrm A}^{\mathrm{iso}}$ region contains a larger contribution from AB linear diblock configurations, with $f^*\approx0$.
% This is because the A-shell-forming geometry already favors lower, or negative, spontaneous curvature, so the system can maintain a low interfacial energy without converting strongly toward A$_2$B-like topologies.
% Consequently, the onset of $f^*>0$ is delayed to around $\phi_{\mathrm A}^{\mathrm{iso}}\approx0.50$, and the variation of $f_\Omega^\ast$ becomes smoother than that for $\bar{\phi}_{\mathrm A}=0.42$, indicating a more stable curvature-driven conformational adjustment.
The response of $P_\Omega(f)$ to $\phi_{\mathrm A}^{\mathrm{iso}}$ further reveals the interfacial-curvature response of L--S TT copolymers.
This response differs between the two O$^{70}$ states because the A component plays distinct structural roles in the A-core-forming and A-shell-forming geometries.
For $\bar{\phi}_{\mathrm A}=0.42$, where A forms the network core, $f_\Omega^\ast$ shifts from the A$_2$B-like endpoint to the AB-like limit near $\phi_{\mathrm A}^{\mathrm{iso}}\approx0.23$, remains close to zero over an intermediate range, and then increases toward an A$_2$AB-like topology at higher $\phi_{\mathrm A}^{\mathrm{iso}}$.
This behavior is consistent with the larger spatial extent of the two A branches in the A$_2$B-like topology, which favors a lower, or more negative, A-side spontaneous curvature.
This lower-curvature tendency can compensate for the more positive curvature required by the A-core geometry.
As the sampled layer moves inward toward the A-rich core, the preferred topology shifts from the AB-like limit toward A$_2$AB-like configurations, whose two A branches contribute more strongly to filling the A-rich region.
By contrast, for $\bar{\phi}_{\mathrm A}=0.55$, where A forms the shell, the low-$\phi_{\mathrm A}^{\mathrm{iso}}$ region retains a larger AB-like contribution with $f_\Omega^\ast\approx0$.
Because the A-shell geometry favors lower, or more negative, A-side curvature, the topology redistribution required to accommodate the local geometry differs from that in the A-core state.
Consequently, the onset of finite $f_\Omega^\ast$ occurs at a higher $\phi_{\mathrm A}^{\mathrm{iso}}$, around $0.50$, and its variation with density level is smoother than in the Core state.
These contrasting responses therefore reflect the different curvature requirements of the A-core and A-shell geometries and the corresponding redistribution among accessible topologies with different curvature preferences.
Because $P_\Omega(f)$ characterizes this layer dependent topology redistribution but does not directly quantify the local interfacial geometry, we further examine the relation between topology selection and local curvature.


To characterize the interfacial geometry directly, we calculate the local mean curvature from the A-component density field.
For an A-component isodensity surface, the unit normal is defined as $\bm n(\bm r)=-\nabla\phi_{\mathrm A}(\bm r)/|\nabla\phi_{\mathrm A}(\bm r)|$, pointing from the A-rich region toward the B-rich region.
The local mean curvature is then evaluated as~\cite{Cheng2024Understand}
\begin{equation}
    H(\bm r) = \frac{1}{2}\nabla\cdot\bm n(\bm r).
\end{equation}
With this convention, positive $H$ corresponds locally to an A-core-like curvature, whereas negative $H$ corresponds to an A-shell-like curvature.
In the numerical evaluation, the spatial gradients and divergence are computed using periodic spectral derivatives, with a weak Gaussian smoothing applied to $\phi_{\mathrm A}(\bm r)$ to suppress grid-scale fluctuations in the curvature.


\begin{figure}[!htb]
    \centering
    \includegraphics[]{./Figure/O70-Curvature.png}
    \caption{
        (a,b) Mean curvature fields $H (R^{-1}_{\mathrm g})$ on the central $z$ section for the Core ($\bar{\phi}_{\mathrm A}=0.42$) and Shell ($\bar{\phi}_{\mathrm A}=0.55$) states; cyan and magenta contours mark $\phi_{\mathrm A}^{\mathrm{iso}}=0.50$ and $0.80$.
        (c,d) $f_{\mathcal R}^{\ast}(\bm r)$ versus $H (R^{-1}_{\mathrm g})$ on the corresponding three-dimensional isodensity surfaces at $\phi_{\mathrm A}^{\mathrm{iso}}=0.50$ and $0.80$, respectively.
        Dark red and orange symbols denote the Core ($\bar{\phi}_{\mathrm A}=0.42$) and Shell ($\bar{\phi}_{\mathrm A}=0.55$) states, respectively.
    }
    \label{fig:O70Curvature}
\end{figure}


Fig.~\ref{fig:O70Curvature}a,b shows the local mean-curvature fields together with the $\phi_{\mathrm A}^{\mathrm{iso}}=0.50$ and $0.80$ isodensity contours.
The former samples the A/B interfacial region, whereas the latter probes a more A-rich region.
For Fig.~\ref{fig:O70Curvature}c,d, spatial points $\bm r$ are sampled over the entire three-dimensional isodensity surface to obtain $H(\bm r)$ and $P_{\mathcal R(\bm r)}(f)$, from which the corresponding $f_{\mathcal R}^{\ast}(\bm r)$ values are determined.
A clear distinction between the Core and Shell states is observed at the two density levels.
On the interfacial surface, $\phi_{\mathrm A}^{\mathrm{iso}}=0.50$, the Core state is biased toward more positive $H$ and exhibits larger $f_{\mathcal R}^{\ast}$, whereas the Shell state is biased toward more negative $H$ and smaller $f_{\mathcal R}^{\ast}$.
On the more A-rich surface, $\phi_{\mathrm A}^{\mathrm{iso}}=0.80$, the relative ordering of $f_{\mathcal R}^{\ast}$ between the two states is reversed, highlighting the distinct topology responses associated with the A-core and A-shell geometries.
Despite these differences, $f_{\mathcal R}^{\ast}$ shows a weak overall tendency to increase with increasing $H$ within the sampled surfaces, although appreciable scatter remains and the trend becomes less distinct near $H\approx0$.
These results show that the topology--curvature response is strongly conditioned by whether A forms the network core or the surrounding shell, while retaining an overall sensitivity of the preferred topology to the local mean curvature.



Taken together, the asymmetric displacement of the O$^{70}$ phase region reflects the coupling between sliding topology redistribution and the local curvature environment.
When A forms the core, topology selection responds to a more positive curvature environment, whereas when A forms the shell, it is redistributed to accommodate a lower, or more negative, curvature requirement.
The analysis of $H$ and $f_{\mathcal R}^{\ast}$ further shows that the two geometries exhibit distinct responses of topology to curvature, while the preferred topology retains a weaker overall dependence on the local mean curvature.
At larger $\phi_{\mathrm a}=0.75$, topology locking reduces the available curvature adjustment, and the high-$\bar{\phi}_{\mathrm A}$ O$^{70}$ region correspondingly shifts back toward higher $N\chi$ (Fig.~\ref{fig:PhaseDiagram}c).
